\documentclass[openright, oneside, a4paper, article, brazil]{abntex2}
\usepackage{formularios}
% --------- %
% Metadados %
% --------- %
\hypersetup{%
	colorlinks=false,
	allbordercolors=white,
	pdftitle = {Representação},
	pdfauthor = {Conselho Regional de Psicologia de São Paulo},
	pdfsubject = {Representação},
	pdfkeywords = {%
		Representação,
		Código de Ética
	}%
}
\title{Representação	}
\author{Conselho Regional de Psicologia de São Paulo}
% --------- %
% DOCUMENTO %
% --------- %
\begin{document}
	\pagestyle{crp}
	\centering
	\Form

\section{REPRESENTAÇÃO}
\vspace*{\baselineskip}

\setasuspacing{\DoubleSpacing}
\justifying

\preenctext{14}{nome}{47.3em}{Eu, } \hfill ,\\
abaixo assinada/o, venho apresentar, nos termos do art.~59 da Resolução CFP nº~011/2019, \textbf{representação} contra a/o psicóloga/o ou pessoa jurídica abaixo qualificada/o, por possível violação ao Código de Ética Profissional da/o Psicóloga/o.

\begin{framed}
	\noindent
	\begin{minipage}{\linewidth}
		\noindent REPRESENTANTE
		\vspace*{0.5\baselineskip}

		\begin{minipage}{\linewidth}
			\preenctext{14}{nome2}{39.5em}{Nome completo: \hfill}\\
			\preenctext{14}{profissao}{42em}{Profissão: \hfill}\\
			\preenctext{14}{rg}{10em}{RG: } \hfill%
			\preenctext{14}{cpf}{10em}{CPF: } \hfill%
			Data de nascimento: \dataabrev{14}{1}\\
			\preenctext{14}{endereco}{42.5em}{Endereço: \hfill}\\
			\preenctext{14}{complemento}{4em}{Complemento: } \hfill%
			\preenctext{14}{bairro}{18em}{Bairro: } \hfill%
			\preenctext{14}{cep}{8em}{CEP: }\\
			\preenctext{14}{municipio}{36em}{Município: } \hfill%
			\preenctext{14}{uf}{2em}{UF: }\\
			Telefones: \hfill
			\preenctext{14}{telefone_fixo}{10.5em}{fixo: } \hfill%
			\preenctext{14}{celular}{10.5em}{celular: } \hfill%
			\preenctext{14}{profissao}{10.5em}{outro: }\\
			\preenctext{14}{e-mail}{44em}{\textit{e-mail}: \hfill}
		\end{minipage}
	\end{minipage}
\end{framed}

\begin{framed}%
	\noindent
	\begin{minipage}{\linewidth}
		\noindent REPRESENTADA/O
		\vspace*{0.5\baselineskip}

			\begin{minipage}{\linewidth}
		\preenctext{14}{nome2}{39.5em}{Nome completo: \hfill}\\
		\preenctext{14}{profissao}{10em}{Inscrição no CRP-06: }\\
		\preenctext{14}{endereco}{37.5em}{Endereço (se souber): \hfill}\\
		\preenctext{14}{complemento}{4em}{Complemento: } \hfill%
		\preenctext{14}{bairro}{18em}{Bairro: } \hfill%
		\preenctext{14}{cep}{8em}{CEP: }\\
		\preenctext{14}{municipio}{36em}{Município: } \hfill%
		\preenctext{14}{uf}{2em}{UF: }\\
Telefones: \hfill
\preenctext{14}{telefone_fixo}{10.5em}{fixo: } \hfill%
\preenctext{14}{celular}{10.5em}{celular: } \hfill%
\preenctext{14}{profissao}{10.5em}{outro: }\\
		\preenctext{14}{e-mail}{44em}{\textit{e-mail}: \hfill}
	\end{minipage}
	\end{minipage}
\end{framed}

Declaro, para fins de direito, que as informações e documentos enviados na representação são verdadeiros e autênticos.

Li e estou ciente de que a representação no Conselho Regional de Psicologia tramita de acordo com o \textbf{Código de Processamento Disciplinar} (CPD; Resolução CFP nº~11/2019), disponível no site www.crpsp.org.br --- item Orientação/legislação.

\subsection{Em anexo}

\begin{itemize}
	\item Descrição detalhada dos fatos envolvendo o exercício profissional da/o psicóloga/o	(obrigatória);
	\item prova documental que possa servir à apuração do(s) fato(s) e de sua autoria;
	\item rol (relação) de testemunhas (se houver);
	\item manifestar interesse em participar de mediação com a/o representada/o (se houver).
\end{itemize}
\vspace{2\baselineskip}

\dataextenso{1}\\\vspace{-1ex}
\small\hspace{8.6em}cidade\hspace{11em}dia\hspace{7em}mês\hspace{7.15em}ano
\vspace{3.5\baselineskip}

\rule{20em}{0.4pt}\\
\small Assinatura da/o representante

\end{document}